% D2-D9-D1-S12Agent
% Auto-generated from markdown source
% origin: ai

\documentclass[12pt, a4paper]{article}
\usepackage[utf8]{inputenc}
\usepackage[T1]{fontenc}
\usepackage{amsmath, amssymb, amsfonts, amsthm}
\usepackage{textcomp}
\usepackage{geometry}
\usepackage{hyperref}
\usepackage{booktabs}
\usepackage{graphicx}
\usepackage{fancyhdr}
\usepackage{tcolorbox}
\usepackage{enumitem}

\geometry{margin=1in}
\pagestyle{fancy}
\fancyhf{}
\rhead{D2-D9-D1-S12Agent}
\lhead{\thepage}

\newtheorem{definition}{Definition}[section]
\newtheorem{theorem}{Theorem}[section]
\newtheorem{corollary}{Corollary}[section]
\newtheorem{proposition}{Proposition}[section]
\newtheorem{lemma}{Lemma}[section]
\newtheorem{conjecture}{Conjecture}[section]

\newcommand{\calI}{\mathcal{I}}
\newcommand{\calP}{\mathcal{P}}
\newcommand{\calE}{\mathcal{E}}
\newcommand{\calS}{\mathcal{S}}
\newcommand{\calH}{\mathcal{H}}
\newcommand{\calO}{\mathcal{O}}
\newcommand{\calF}{\mathcal{F}}
\newcommand{\calU}{\mathcal{U}}
\newcommand{\calV}{\mathcal{V}}
\newcommand{\calA}{\mathcal{A}}
\newcommand{\calG}{\mathcal{G}}
\newcommand{\calB}{\mathcal{B}}
\newcommand{\calC}{\mathcal{C}}
\newcommand{\calL}{\mathcal{L}}
\newcommand{\calM}{\mathcal{M}}
\newcommand{\calX}{\mathcal{X}}
\newcommand{\calY}{\mathcal{Y}}
\newcommand{\calW}{\mathcal{W}}
\newcommand{\calT}{\mathcal{T}}
\newcommand{\calZ}{\mathcal{Z}}
\newcommand{\calR}{\mathcal{R}}
\newcommand{\calQ}{\mathcal{Q}}
\newcommand{\calD}{\mathcal{D}}
\newcommand{\calN}{\mathcal{N}}
\newcommand{\calK}{\mathcal{K}}

\newcommand{\Tr}{\text{Tr}}
\newcommand{\E}{\mathbb{E}}
\newcommand{\diff}{\mathrm{d}}
\newcommand{\KL}{D_{\mathrm{KL}}}
\newcommand{\argmax}{\arg\max}
\newcommand{\argmin}{\arg\min}
\newcommand{\hbarF}{\hbar}

\newtcolorbox{thmbox}[1]{
  colback=blue!5!white, colframe=blue!75!black, fonttitle=\bfseries,
  title={#1}, breakable
}
\newtcolorbox{defbox}[1]{
  colback=green!5!white, colframe=green!75!black, fonttitle=\bfseries,
  title={#1}, breakable
}

\begin{document}

\title{D2-D9-D1-S12Agent}
\author{RALPH Loop}
\date{}
\maketitle

\begin{abstract}
This paper presents D2-D9-D1-S12Agent, formalizing the underlying principles using information-theoretic methods.

\textbf{Keywords}: D2-D9-D1-S12Agent, information theory, formalization.
\end{abstract}

\tableofcontents
\section{D2-D9-D1-S12Agent}


****D2-D9-D1-S12-r01joint-agent-emergence-framework


\#\#

\textbf{D2-1/D2-2}  Agent \textbf{ + }
\textbf{D9-IV} ****
\textbf{D1-S12}  chimera **** =

\textbf{*}\textbf{*}**


\#\#

\#\#\#

\textbf{ Agent A  \(\mathcal{R} = f_\theta(\mathcal{X})\) \(|A|\) \(t\)}

\[\text{Intelligence}(A) \equiv \bigwedge_{i=1}^{4} C_i(A)\]

****

|  |  |  |  |
|------|------|------|----------|
| \textbf{C1} | Agent  | D2-1 | \(|A| > I_{plan}^* + \delta_{margin}\) |
| \textbf{C2} |  | D2-2 | \(t > t_c\) |
| \textbf{C3} |  | D9-IV | \(S(f_\theta) = (\mathcal{D}, \mathcal{L}, \mathcal{T}) > \Sigma_c\) |
| \textbf{C4} |  | D1-S12 | \(N \geq N_c \approx 500\) |


\#\#

\subsubsection{1. -D9-IV + D1-S12}


\textbf{*}*D1-S12-T1  D9-IV  \(\mathcal{T}\) **

\textbf{D9-IV } \(\mathcal{T}\)

\textbf{D1-S12-T1 } =  \textbf{chimera-like spatially-heterogeneous sync states}

****
\[\mathcal{T}(A) \equiv \sup_{K(x)} P_{\text{chimera}}(A \mid K(x), N, \epsilon, T_{sim})\]

\(N < N_c\) \(\mathcal{T}(A) = 0\)

\subsubsection{2. -D2-1 + D9-IV}


****-
\[\text{ } \mathcal{D}(f_\theta) < \mathcal{D}_c \implies \text{ } |A| > |A|_c \text{Agent }\]

****
1.  \(\mathcal{D}\) D9-IV  1
2.
3. **** =

\subsubsection{3. -D2-2 + D1-S12}


\textbf{** \(t_c\) } \(\mathcal{T}\)**

****
\[t_c(\mathcal{T}) = \frac{1}{\lambda} \log\left(\frac{K^* - K_0}{K^* - K_c(\mathcal{T})}\right)\]

\(\mathcal{T}\)  chimera-like \(K_c(\mathcal{T}) < K_c(\text{uniform})\)  ****

\subsubsection{4.}


****
\[\Psi_{\text{intelligence}} = \prod_{i=1}^{4} \Phi\left(\frac{X_i - X_i^c}{X_i^c}\right)\]

\(X_1 = |A|\[X_2 = t\]X_3 = \Sigma(f_\theta)\)\(X_4 = N\)

\textbf{*}***


\#\#

\#\#\#

\textbf{ 3 }

****
1.  3
2. \textbf{ 1} +   D9-IV  3
3. \textbf{ 2} +   D2-2
4. \textbf{ 3} +   D1-S12 gradient_locking
5. ****


\#\#



1. \textbf{D2-1 } \(|A| < |A|_c\)
2. \textbf{D2-2 } \(t_c\)
3. \textbf{D9-IV }\(\epsilon < |\text{Cov}(d_i,d_j)| < \delta\)
4. \textbf{D1-S12 } \(N < 500\)  chimera gradient_locking
5. ****


\subsection{ Agent}


|  |  |  |
|------|----------|----------|
|  |  \(|A| > I_{plan}^* + 10\\%\) |  |
|  |  \(t > t_c\)  |  +  RAG  |
|  |  \(\mathcal{D}\)  |  |
|  |  \(N \geq 500\)  |  |


\#\#

1. ****
2. \textbf{ D9-IV "" chimera }
3. \textbf{ D2-1/D2-2  +  = }
4. \textbf{N}


* D2-1/D2-2 + D9-IV + D1-S12 *
*2026-05-04 | r01 | *

\begin{thebibliography}{99}
\bibitem{ref1} Shannon, C.E. (1948). A Mathematical Theory of Communication.
\bibitem{ref2} Cover, T.M. (2006). Elements of Information Theory.
\end{thebibliography}

\end{document}